% Options for packages loaded elsewhere
\PassOptionsToPackage{unicode}{hyperref}
\PassOptionsToPackage{hyphens}{url}
\PassOptionsToPackage{space}{xeCJK}
\documentclass[
]{article}
\usepackage{xcolor}
\usepackage{amsmath,amssymb}
\setcounter{secnumdepth}{-\maxdimen} % remove section numbering
\usepackage{iftex}
\ifPDFTeX
  \usepackage[T1]{fontenc}
  \usepackage[utf8]{inputenc}
  \usepackage{textcomp} % provide euro and other symbols
\else % if luatex or xetex
  \usepackage{unicode-math} % this also loads fontspec
  \defaultfontfeatures{Scale=MatchLowercase}
  \defaultfontfeatures[\rmfamily]{Ligatures=TeX,Scale=1}
\fi
\usepackage{lmodern}
\ifPDFTeX\else
  % xetex/luatex font selection
  \setmainfont[]{Times New Roman}
  \ifXeTeX
    \usepackage{xeCJK}
    \setCJKmainfont[]{PingFang SC}
  \fi
  \ifLuaTeX
    \usepackage[]{luatexja-fontspec}
    \setmainjfont[]{PingFang SC}
  \fi
\fi
% Use upquote if available, for straight quotes in verbatim environments
\IfFileExists{upquote.sty}{\usepackage{upquote}}{}
\IfFileExists{microtype.sty}{% use microtype if available
  \usepackage[]{microtype}
  \UseMicrotypeSet[protrusion]{basicmath} % disable protrusion for tt fonts
}{}
\makeatletter
\@ifundefined{KOMAClassName}{% if non-KOMA class
  \IfFileExists{parskip.sty}{%
    \usepackage{parskip}
  }{% else
    \setlength{\parindent}{0pt}
    \setlength{\parskip}{6pt plus 2pt minus 1pt}}
}{% if KOMA class
  \KOMAoptions{parskip=half}}
\makeatother
\setlength{\emergencystretch}{3em} % prevent overfull lines
\providecommand{\tightlist}{%
  \setlength{\itemsep}{0pt}\setlength{\parskip}{0pt}}
\usepackage{bookmark}
\IfFileExists{xurl.sty}{\usepackage{xurl}}{} % add URL line breaks if available
\urlstyle{same}
\hypersetup{
  hidelinks,
  pdfcreator={LaTeX via pandoc}}

\author{}
\date{}

\begin{document}

\section{普惠金融风险动态评估平台指标口径与生成原理详细说明（专业公式版）}\label{ux666eux60e0ux91d1ux878dux98ceux9669ux52a8ux6001ux8bc4ux4f30ux5e73ux53f0ux6307ux6807ux53e3ux5f84ux4e0eux751fux6210ux539fux7406ux8be6ux7ec6ux8bf4ux660eux4e13ux4e1aux516cux5f0fux7248}

\subsection{一、文档目的}\label{ux4e00ux6587ux6863ux76eeux7684}

本文档用于系统说明平台中``处理后数据''的生成逻辑，重点回答以下问题：

\begin{enumerate}
\def\labelenumi{\arabic{enumi}.}
\tightlist
\item
  原始数据从何而来；
\item
  原始数据如何被清洗、标准化、筛选与排序；
\item
  首页六大核心指标、雷达图、原始检索卡片、区域画像、趋势预测、政策模拟等结果分别通过什么方法得到；
\item
  当前系统中哪些结果属于真实统计数据直接驱动，哪些属于代理变量估计，哪些属于基于真实样本的建模展示。
\end{enumerate}

因此，本文档属于``方法论与指标口径说明''，适用于答辩材料、技术附录或队内原理解释。

\begin{center}\rule{0.5\linewidth}{0.5pt}\end{center}

\subsection{二、原始数据层与结构化底表}\label{ux4e8cux539fux59cbux6570ux636eux5c42ux4e0eux7ed3ux6784ux5316ux5e95ux8868}

\subsubsection{2.1 数据来源}\label{ux6570ux636eux6765ux6e90}

当前平台真实接入的数据底座以广东省官方公开统计样本为核心，后端数据文件位置为：

\begin{itemize}
\tightlist
\item
  \texttt{assets/data/guangdong\_official\_raw\_data.json}
\end{itemize}

该底表由后端函数 \texttt{\_load\_guangdong\_raw\_data()}
统一读取，并组织为：

\begin{itemize}
\tightlist
\item
  \texttt{metadata}：记录数据集名称、版本号、采集时间、所属区域等元信息；
\item
  \texttt{records}：记录结构化统计样本。
\end{itemize}

\subsubsection{2.2
单条样本记录的字段结构}\label{ux5355ux6761ux6837ux672cux8bb0ux5f55ux7684ux5b57ux6bb5ux7ed3ux6784}

根据后端检索逻辑，单条记录可抽象为：

\[
r_i = \left( p_i, c_i, h_i, m_i, v_i, u_i, t_i, f_i, y_i, g_i, s_i, l_i \right)
\]

其中：

\begin{itemize}
\tightlist
\item
  \(p_i\)：省份（province）
\item
  \(c_i\)：城市（city）
\item
  \(h_i\)：行业（industry）
\item
  \(m_i\)：指标名（indicator）
\item
  \(v_i\)：指标值（value）
\item
  \(u_i\)：单位（unit）
\item
  \(t_i\)：时间口径（period）
\item
  \(f_i\)：频率（frequency）
\item
  \(y_i\)：同比增速（yoy\_pct）
\item
  \(g_i\)：指标类别（category）
\item
  \(s_i\)：来源机构与标题（source）
\item
  \(l_i\)：来源链接（source\_url）
\end{itemize}

因此，平台当前的数据不是截图数字，也不是静态写死文本，而是结构化底表驱动的记录集合：

\[
\mathcal{R} = \{r_1, r_2, \dots, r_n\}
\]

\begin{center}\rule{0.5\linewidth}{0.5pt}\end{center}

\subsection{三、原始数据的清洗、标准化与筛选}\label{ux4e09ux539fux59cbux6570ux636eux7684ux6e05ux6d17ux6807ux51c6ux5316ux4e0eux7b5bux9009}

\subsection{3.1
行政区划标准化}\label{ux884cux653fux533aux5212ux6807ux51c6ux5316}

由于原始公开数据在地区名称层面可能存在不同写法，系统先做名称统一。城市标准化函数可写作：

\[
\phi_{city}(x)=
\begin{cases}
\text{广州市}, & x=\text{广州}\\
\text{深圳市}, & x=\text{深圳}\\
\text{佛山市}, & x=\text{佛山}\\
\text{东莞市}, & x=\text{东莞}\\
\text{珠海市}, & x=\text{珠海}\\
x, & \text{otherwise}
\end{cases}
\]

这一步的作用是保证前端地图点击、筛选框选择与后端查询字段严格一致。

\subsection{3.2
行业主类映射}\label{ux884cux4e1aux4e3bux7c7bux6620ux5c04}

平台并不要求用户直接选择所有细分行业，而是通过主行业选项映射到多个统计子行业。可定义行业映射函数：

\[
\phi_{industry}(H)=\{h_1,h_2,\dots,h_k\}
\]

例如：

\begin{itemize}
\tightlist
\item
  ``制造业''映射到``先进制造业''\,``高技术制造业''\,``食品制造业''\,``汽车制造业''等；
\item
  ``软件和信息技术服务业''映射到``信息传输、软件和信息技术服务业''\,``科学研究和技术服务业''等。
\end{itemize}

于是记录 \(r_i\) 与行业筛选条件 \(H\) 匹配，当且仅当：

\[
h_i \in \phi_{industry}(H)
\]

\subsection{3.3 结构化筛选}\label{ux7ed3ux6784ux5316ux7b5bux9009}

对于给定筛选条件：

\[
F=(P,C,H,Y)
\]

其中：

\begin{itemize}
\tightlist
\item
  \(P\)：省份
\item
  \(C\)：城市
\item
  \(H\)：行业
\item
  \(Y\)：年份
\end{itemize}

平台从原始记录集 \(\mathcal{R}\) 中构造筛选后的记录子集：

\[
\mathcal{R}_F = \left\{ r_i \in \mathcal{R} \mid p_i=P,\ c_i=C,\ h_i \in \phi_{industry}(H),\ t_i \in Y \right\}
\]

若年份选择为``全部年份''，则不约束
\(t_i\)；若为``2023及以前''，则保留所有不属于 2024、2025、2026 的记录。

\subsection{3.4
检索结果优先级排序}\label{ux68c0ux7d22ux7ed3ux679cux4f18ux5148ux7ea7ux6392ux5e8f}

原始数据检索模块并非简单返回全部匹配记录，而是进行评分排序。对于记录
\(r_i\)，定义总分：

\[
S_i = S_i^{match} + S_i^{priority}
\]

其中：

\begin{itemize}
\tightlist
\item
  \(S_i^{match}\)：关键词匹配得分
\item
  \(S_i^{priority}\)：业务优先级得分
\end{itemize}

优先级加分规则体现为：

\begin{itemize}
\tightlist
\item
  地区生产总值权重最高；
\item
  增加值、社会消费品零售总额、营业收入、流动资产合计、用电量、存贷款余额等依次加权；
\item
  主行业、城市级记录获得额外权重。
\end{itemize}

若记录的值满足：

\[
v_i \in \{\varnothing, "", \text{未披露}\}
\]

则该记录直接剔除，不进入前端展示结果集。

\begin{center}\rule{0.5\linewidth}{0.5pt}\end{center}

\subsection{四、首页六个核心指标的生成原理}\label{ux56dbux9996ux9875ux516dux4e2aux6838ux5fc3ux6307ux6807ux7684ux751fux6210ux539fux7406}

首页六个核心指标由 \texttt{\_build\_real\_dashboard\_overview()}
在用户筛选后实时重算。其核心思想是：从 \(\mathcal{R}_F\)
中抽取最能代表当前场景的关键统计指标，再通过代理变量映射和线性合成得到标准化得分。

\subsubsection{4.1
底层关键记录抽取}\label{ux5e95ux5c42ux5173ux952eux8bb0ux5f55ux62bdux53d6}

平台从筛选记录中选择如下关键量：

\begin{itemize}
\tightlist
\item
  地区生产总值：\(GDP\)
\item
  社会消费品零售总额：\(RET\)
\item
  存款余额：\(DEP\)
\item
  贷款余额：\(LOAN\)
\item
  营业收入：\(REV\)
\item
  流动资产合计：\(CA\)
\item
  用电量：\(ELEC\)
\item
  规模以上工业销售产值：\(SALE\)
\item
  增加值：\(AV\)
\end{itemize}

对应同比增速记为：

\[
g, r, d, l, v, a, e, s, q
\]

分别对应：

\begin{itemize}
\tightlist
\item
  \(g\)：GDP 同比
\item
  \(r\)：社零同比
\item
  \(d\)：存款同比
\item
  \(l\)：贷款同比
\item
  \(v\)：营业收入同比
\item
  \(a\)：流动资产同比
\item
  \(e\)：用电量同比
\item
  \(s\)：销售产值同比
\item
  \(q\)：增加值同比
\end{itemize}

\begin{center}\rule{0.5\linewidth}{0.5pt}\end{center}

\subsubsection{4.2
区域经济景气度}\label{ux533aux57dfux7ecfux6d4eux666fux6c14ux5ea6}

定义区域经济景气度为：

\[
ECO = \operatorname{clip}\left(72 + 3.2g + 1.6r,\ 0,\ 100\right)
\]

解释：

\begin{itemize}
\tightlist
\item
  基础分为 72；
\item
  GDP 同比是主导项，因此系数较高；
\item
  社零同比反映需求侧修复，因此作为辅助项。
\end{itemize}

该指标本质上是以宏观产出与消费恢复共同表征区域景气水平。

\begin{center}\rule{0.5\linewidth}{0.5pt}\end{center}

\subsubsection{4.3
行业发展健康指数}\label{ux884cux4e1aux53d1ux5c55ux5065ux5eb7ux6307ux6570}

定义行业健康指数为：

\[
IND = \operatorname{clip}\left(68 + 1.7s + 1.2v + 0.9e + 1.1q,\ 0,\ 100\right)
\]

解释：

\begin{itemize}
\tightlist
\item
  销售产值同比 \(s\) 是最重要信号；
\item
  营业收入同比 \(v\) 反映企业经营恢复；
\item
  用电量同比 \(e\) 作为生产活跃度代理；
\item
  增加值同比 \(q\) 体现增长质量。
\end{itemize}

这是一种``产出-营收-活跃度-附加值''四维联合代理。

\begin{center}\rule{0.5\linewidth}{0.5pt}\end{center}

\subsubsection{4.4
普惠金融环境指数}\label{ux666eux60e0ux91d1ux878dux73afux5883ux6307ux6570}

定义金融环境指数为：

\[
FIN = \operatorname{clip}\left(70 + 1.3d + 1.0l + 0.8\max(0,r),\ 0,\ 100\right)
\]

解释：

\begin{itemize}
\tightlist
\item
  存款同比 \(d\) 反映区域资金沉淀能力；
\item
  贷款同比 \(l\) 反映信贷投放能力；
\item
  正向消费修复 \(\max(0,r)\) 代表资金流转活跃度。
\end{itemize}

因此，该指标不是某个官方``金融环境现成分数''，而是通过金融供给与消费修复构造出的代理指标。

\begin{center}\rule{0.5\linewidth}{0.5pt}\end{center}

\subsubsection{4.5
小微企业经营健康}\label{ux5c0fux5faeux4f01ux4e1aux7ecfux8425ux5065ux5eb7}

定义小微经营健康指数为：

\[
MIC = \operatorname{clip}\left(64 + 1.4v + 0.8a - 0.6\max(0,-s),\ 0,\ 100\right)
\]

解释：

\begin{itemize}
\tightlist
\item
  营业收入同比 \(v\) 上升，说明经营恢复；
\item
  流动资产同比 \(a\) 上升，说明短期周转改善；
\item
  若销售产值同比为负，则说明订单端走弱，因此通过 \(\max(0,-s)\) 扣分。
\end{itemize}

该指标反映的是小微企业经营活跃度与现金流稳定性的综合状态。

\begin{center}\rule{0.5\linewidth}{0.5pt}\end{center}

\subsubsection{4.6
普惠小微贷款不良率}\label{ux666eux60e0ux5c0fux5faeux8d37ux6b3eux4e0dux826fux7387}

在当前筛选态下，不良率不是直接从底表中读取，而是通过金融环境与经营健康代理测算得到：

\[
NPL = \operatorname{clip}\left(3.15 - 0.011FIN - 0.008MIC + 0.03\max(0,-v),\ 0.9,\ 3.8\right)
\]

解释：

\begin{itemize}
\tightlist
\item
  金融环境越好，不良率越低；
\item
  小微经营越健康，不良率越低；
\item
  若营业收入同比为负，则经营承压，不良率上升。
\end{itemize}

因此该值属于``真实统计样本驱动下的代理测算结果''，而不是一条原始公报字段的直接拷贝。

\begin{center}\rule{0.5\linewidth}{0.5pt}\end{center}

\subsubsection{4.7
政策实施有效性}\label{ux653fux7b56ux5b9eux65bdux6709ux6548ux6027}

政策实施有效性定义为多维综合评分：

\[
POL = \operatorname{clip}\left(0.22ECO + 0.26IND + 0.22FIN + 0.18MIC + 0.12(4-NPL)\times 18,\ 0,\ 100\right)
\]

解释：

\begin{itemize}
\tightlist
\item
  区域景气、行业健康、金融环境、小微经营共同决定政策作用空间；
\item
  不良率经反向转换后纳入综合评价；
\item
  行业健康权重最高，说明平台认为政策效果首先体现为行业修复与经营改善。
\end{itemize}

\begin{center}\rule{0.5\linewidth}{0.5pt}\end{center}

\subsubsection{4.8
指标趋势值与方向}\label{ux6307ux6807ux8d8bux52bfux503cux4e0eux65b9ux5411}

趋势箭头不是人工指定，而是由同比值自动生成：

\[
T(x) =
\begin{cases}
(\text{up}, +|x|), & x \ge 0\\
(\text{down}, -|x|), & x < 0
\end{cases}
\]

不良率单独采用``下降为好''的口径，因此趋势文字保持负值显示，方向固定为下降。

\begin{center}\rule{0.5\linewidth}{0.5pt}\end{center}

\subsection{五、雷达图是如何映射出来的}\label{ux4e94ux96f7ux8fbeux56feux662fux5982ux4f55ux6620ux5c04ux51faux6765ux7684}

雷达图不是重新训练一套指标，而是对首页指标做二次映射。其六个维度为：

\[
RADAR = \{ECO, IND, FIN, MIC, CRD, POL\}
\]

其中信用环境指数 \(CRD\) 定义为：

\[
CRD = 0.58FIN + 0.42MIC
\]

解释：

平台认为``信用环境''不仅取决于金融供给，也取决于企业自身经营健康，因此用金融环境与小微经营做加权融合。

\begin{center}\rule{0.5\linewidth}{0.5pt}\end{center}

\subsection{六、原始数据卡片是如何生成的}\label{ux516dux539fux59cbux6570ux636eux5361ux7247ux662fux5982ux4f55ux751fux6210ux7684}

前端原始数据检索卡片不是二次计算结果，而是结构化记录的直接渲染。对于单条记录
\(r_i\)，卡片展示内容对应为：

\begin{itemize}
\tightlist
\item
  标题：\(h_i + m_i\)
\item
  主值：\(v_i + u_i\)
\item
  时间：\(t_i\)
\item
  标签：\(g_i, f_i, y_i\)
\item
  说明：\(s_i\)
\item
  来源链接：\(l_i\)
\end{itemize}

因此原始数据检索模块本质是：

\[
\mathcal{R}_F \xrightarrow{\text{排序与过滤}} \mathcal{R}_F^{top} \xrightarrow{\text{卡片渲染}} UI
\]

\begin{center}\rule{0.5\linewidth}{0.5pt}\end{center}

\subsection{七、地图、区域侧栏与区域画像的生成逻辑}\label{ux4e03ux5730ux56feux533aux57dfux4fa7ux680fux4e0eux533aux57dfux753bux50cfux7684ux751fux6210ux903bux8f91}

\subsubsection{7.1 全国热力图}\label{ux5168ux56fdux70edux529bux56fe}

全国热力图由 \texttt{get\_risk\_map()} 提供。当前热力值为省级样本分值：

\[
MAP = \{(Beijing,85), (Shanghai,92), (Guangdong,95), \dots\}
\]

这部分当前主要属于``区域建模展示层''，并非全国全量真实底表实时聚合结果。其作用是表达区域差异与空间风险分布。

\subsubsection{7.2
地图脉冲散点}\label{ux5730ux56feux8109ux51b2ux6563ux70b9}

地图上的脉冲散点可表示为：

\[
A_j = (lon_j, lat_j, w_j, desc_j)
\]

其中：

\begin{itemize}
\tightlist
\item
  \(lon_j, lat_j\)：地理坐标
\item
  \(w_j\)：脉冲权重
\item
  \(desc_j\)：异常说明
\end{itemize}

该模块用于表现``长期环境较优但短期链路波动上升''的告警概念。

\subsubsection{7.3
省份侧栏画像}\label{ux7701ux4efdux4fa7ux680fux753bux50cf}

点击地图后右侧侧栏调用
\texttt{get\_province\_detail()}。这部分数据由三块样本组成：

\begin{itemize}
\tightlist
\item
  \texttt{\_PROVINCE\_DATA}
\item
  \texttt{\_INDUSTRY\_PIE}
\item
  \texttt{\_RISK\_FACTORS}
\end{itemize}

例如广东对应的核心值为：

\[
(NPL, CRD, MIC) = (1.85,\ 95,\ 82)
\]

侧栏饼图和柱图则分别映射行业结构分布与风险因子贡献。

需要说明的是，这部分属于``省份画像样本 + 规则映射''的区域画像展示层。

\begin{center}\rule{0.5\linewidth}{0.5pt}\end{center}

\subsection{八、不良率趋势预测与 LSTM
模块}\label{ux516bux4e0dux826fux7387ux8d8bux52bfux9884ux6d4bux4e0e-lstm-ux6a21ux5757}

平台趋势图接口为 \texttt{get\_trends()}。当前不良率历史样本序列为：

\[
\mathbf{x} = [2.50, 2.45, 2.40, 2.38, 2.35, 2.30, 2.25, 2.28, 2.20, 2.18, 2.15, 2.14]
\]

预测序列为：

\[
\hat{\mathbf{x}} = [2.15, 2.14, 2.12, 2.05, 1.95]
\]

并同时给出上下置信区间：

\[
CI^{upper},\ CI^{lower}
\]

系统在接口中明确标注：

\begin{itemize}
\tightlist
\item
  模型：\texttt{LSTM\ (128\ units,\ 2\ layers,\ DropOut=0.2)}
\item
  置信度：\texttt{95.8\%}
\end{itemize}

因此，趋势图属于``演示型历史样本 + LSTM
预测结果''的组合模块，重点在于表达平台具备时间序列预测能力。

\begin{center}\rule{0.5\linewidth}{0.5pt}\end{center}

\subsection{九、政策模拟器的详细生成原理}\label{ux4e5dux653fux7b56ux6a21ux62dfux5668ux7684ux8be6ux7ec6ux751fux6210ux539fux7406}

政策模拟模块位于
\texttt{\_build\_policy\_result()}，是平台中公式链条最完整的部分。

\subsubsection{9.1 输入变量}\label{ux8f93ux5165ux53d8ux91cf}

用户在前端调整三个输入：

\begin{itemize}
\tightlist
\item
  政府担保拨备率 \(g\)
\item
  利率优惠幅度 \(r\)
\item
  补贴等级 \(s\)
\end{itemize}

其中当前系统基准值分别围绕：

\begin{itemize}
\tightlist
\item
  \(g_0 = 10\%\)
\item
  \(r_0 = -0.5\%\)
\item
  \(s_0 = 1\)
\end{itemize}

\subsubsection{9.2 PSM-DID
逻辑下的净政策效应}\label{psm-did-ux903bux8f91ux4e0bux7684ux51c0ux653fux7b56ux6548ux5e94}

当前系统用一个简化净效应函数表示 DID 结果：

\[
\Delta_{did} = 0.018(g-10) + 0.42|r| + 0.12s
\]

解释：

\begin{itemize}
\tightlist
\item
  担保率每高于 10 个百分点，净效应线性上升；
\item
  利率优惠幅度越大，融资缓释作用越强；
\item
  补贴等级越高，产业恢复作用越明显。
\end{itemize}

这一表达并非严格的学术计量回归估计式，而是对 PSM-DID
逻辑的简化政策推演函数，用于在竞赛环境中实现可解释的变量联动。

\subsubsection{9.3
政策实施后的不良率}\label{ux653fux7b56ux5b9eux65bdux540eux7684ux4e0dux826fux7387}

设当前基准不良率为：

\[
NPL_0 = 2.14
\]

则政策实施后预测不良率为：

\[
NPL_1 = \max(0.8,\ NPL_0 - \Delta_{did})
\]

这里加入下界
0.8，是为了保证预测结果符合金融常识，不出现不合理的超低不良率。

\subsubsection{9.4
财政成本函数}\label{ux8d22ux653fux6210ux672cux51fdux6570}

财政成本定义为：

\[
C = 0.8g + 2.5s
\]

单位为亿元。

解释：

\begin{itemize}
\tightlist
\item
  担保拨备率带来财政风险准备支出；
\item
  补贴等级带来财政补贴强度支出。
\end{itemize}

\subsubsection{9.5
受益企业覆盖面}\label{ux53d7ux76caux4f01ux4e1aux8986ux76d6ux9762}

受益企业数定义为：

\[
B = 12000 + 8000\Delta_{did}
\]

并取整数：

\[
B^{*} = \lfloor B \rfloor
\]

解释：

政策净效应越大，理论上可覆盖的企业越多。

\subsubsection{9.6 ROI 口径}\label{roi-ux53e3ux5f84}

平台当前统一 ROI 口径为``每亿元财政投入收益企业数''，定义为：

\[
ROI = \frac{B^{*}}{C + 0.1}
\]

其中 \(0.1\) 是为了避免极端小成本场景下分母过小导致结果不稳定。

这一定义强调的是财政投入效率，而非传统金融收益率。

\subsubsection{9.7
三个月政策路径预测}\label{ux4e09ux4e2aux6708ux653fux7b56ux8defux5f84ux9884ux6d4b}

为了在图上表达政策生效后的平滑变化路径，系统定义：

\[
\delta = \frac{NPL_0 - NPL_1}{3}
\]

则未来三期预测路径为：

\[
\hat{NPL}_1 = NPL_0 - \delta
\]

\[
\hat{NPL}_2 = NPL_0 - 2\delta
\]

\[
\hat{NPL}_3 = NPL_1
\]

这使得趋势图不会从当前值直接跳到最终结果，而是呈现出逐步下行轨迹。

\subsubsection{9.8
风险预警函数}\label{ux98ceux9669ux9884ux8b66ux51fdux6570}

系统还定义了若干风险触发规则：

\begin{itemize}
\tightlist
\item
  若 \(g > 20\)，提示地方财政流动性压力；
\item
  若 \(|r| > 1.0\)，提示银行净息差压缩；
\item
  若 \(s \ge 3\)，提示需与专项债协同。
\end{itemize}

因此平台的政策模拟不是只给正面结果，也会提示高强度政策的副作用。

\subsubsection{9.9 评分卡体系}\label{ux8bc4ux5206ux5361ux4f53ux7cfb}

政策评分卡包含四个分数。

\paragraph{效率分}\label{ux6548ux7387ux5206}

\[
Score_{eff} = \min\left(100,\ \frac{ROI}{200}\right)
\]

\paragraph{覆盖分}\label{ux8986ux76d6ux5206}

\[
Score_{cov} = \min\left(100,\ \frac{B^{*}}{220}\right)
\]

\paragraph{实施分}\label{ux5b9eux65bdux5206}

\[
Score_{impl} = \max\left(55,\ 92 - 1.5\max(0,g-18) - 4\max(0,s-2)\right)
\]

解释：

担保率过高、补贴等级过高，都会提高执行难度，因此实施分会下降。

\paragraph{综合分}\label{ux7efcux5408ux5206}

\[
Score_{total} = 0.4Score_{eff} + 0.35Score_{cov} + 0.25Score_{impl}
\]

这意味着平台在政策方案评价时采用的是``效率优先、覆盖次之、实施可行性约束''的综合逻辑。

\begin{center}\rule{0.5\linewidth}{0.5pt}\end{center}

\subsection{十、行业演化视图与产业链图谱的口径说明}\label{ux5341ux884cux4e1aux6f14ux5316ux89c6ux56feux4e0eux4ea7ux4e1aux94feux56feux8c31ux7684ux53e3ux5f84ux8bf4ux660e}

行业演化视图和产业链图谱目前更适合定义为``真实样本驱动的风险建模展示层''。

\subsubsection{10.1
行业演化视图}\label{ux884cux4e1aux6f14ux5316ux89c6ux56fe}

行业演化视图中的：

\begin{itemize}
\tightlist
\item
  风险等级
\item
  风险分值
\item
  风险驱动热力图
\item
  节点热度排序
\item
  结论与政策建议
\end{itemize}

当前主要来自行业场景样本和风险逻辑配置，其作用是展示行业风险沿时间与链路的变化过程，而非逐字段对应统计局原表。

\subsubsection{10.2 产业链图谱}\label{ux4ea7ux4e1aux94feux56feux8c31}

在 \texttt{get\_industry\_graph()} 中，每个节点定义了：

\[
node_i = (risk\_level_i,\ risk\_score_i,\ description_i,\ policy\_focus_i,\ adjustment_i)
\]

其中 \texttt{adjustment\_i} 进一步映射为政策模拟器的参数增量：

\begin{itemize}
\tightlist
\item
  \texttt{guarantee\_rate\_delta}
\item
  \texttt{interest\_offset\_delta}
\item
  \texttt{subsidy\_level\_delta}
\end{itemize}

因此图谱模块本质是``链路画像 +
策略联动''层，不是原始统计底表直接可得结果。

\begin{center}\rule{0.5\linewidth}{0.5pt}\end{center}

\subsection{十一、结果类型边界说明}\label{ux5341ux4e00ux7ed3ux679cux7c7bux578bux8fb9ux754cux8bf4ux660e}

为了保持表述真实，平台中的处理后数据可以分为三类。

\subsubsection{11.1
真实数据直接驱动}\label{ux771fux5b9eux6570ux636eux76f4ux63a5ux9a71ux52a8}

\begin{itemize}
\tightlist
\item
  原始数据检索结果卡片
\item
  首页筛选后的底层样本记录
\item
  参与六大指标计算的
  GDP、社零、营业收入、流动资产、存贷款、用电量等原始统计样本
\end{itemize}

\subsubsection{11.2
真实数据代理测算}\label{ux771fux5b9eux6570ux636eux4ee3ux7406ux6d4bux7b97}

\begin{itemize}
\tightlist
\item
  首页六大核心指标
\item
  雷达图
\item
  筛选态下的不良率
\item
  筛选态下的政策实施有效性
\end{itemize}

\subsubsection{11.3
基于真实样本的建模展示}\label{ux57faux4e8eux771fux5b9eux6837ux672cux7684ux5efaux6a21ux5c55ux793a}

\begin{itemize}
\tightlist
\item
  全国热力图
\item
  地图脉冲预警
\item
  区域画像侧栏
\item
  行业演化视图
\item
  产业链传导图谱
\item
  LSTM 趋势预测展示
\item
  PSM-DID 简化政策模拟器
\end{itemize}

\begin{center}\rule{0.5\linewidth}{0.5pt}\end{center}

\subsection{十二、总结性表达}\label{ux5341ux4e8cux603bux7ed3ux6027ux8868ux8fbe}

若将整个平台的指标生成逻辑概括为一条链路，则可以写成：

\[
\text{官方原始样本}
\rightarrow
\text{字段标准化}
\rightarrow
\text{结构化筛选}
\rightarrow
\text{关键指标抽取}
\rightarrow
\text{代理变量映射}
\rightarrow
\text{综合评分合成}
\rightarrow
\text{风险识别与情景推演}
\]

因此，平台中的处理后结果并不是凭空生成，而是沿着``真实样本
-\textgreater{} 结构化处理 -\textgreater{} 数学映射 -\textgreater{}
决策输出''的链路逐层得到。当前系统已经真实完成广东省样本接入和首页指标重算，并在此基础上实现了区域画像、行业演化、政策模拟和报告导出的完整闭环。

\end{document}
